% -------------------------------------------------------
% ASSIGNED TO: \textcolor{green}{Person 1}
% -------------------------------------------------------
% INSTRUCTION: Keep existing content. Add scope and limitations
% paragraph at the end (marked in red below).
% Convert first-person to third person throughout.
% -------------------------------------------------------

Currently, there is no standardised diagnostic framework that uses a model's generalisation gap
as a predictor of explanation quality. While regularisation is known to improve accuracy, research
exploring how optimisers like Sharpness-Aware Minimisation (SAM) can be used to smooth the
explanation landscape to produce rules that are faithful and comprehensible is limited
\cite{Dombrowski2022, Springer2024}. Without a ``Red XAI'' audit linking model health to rule quality,
clinicians risk relying on explanations that may be mathematically convenient but biologically
fraudulent \cite{Longo2024}.

The primary research question is as follows: to what extent can the generalisation gap serve as a
diagnostic indicator of explanation fragility, and can SAM be used to improve the fidelity and
comprehensibility of rules extracted using evolutionary algorithms?

This breaks down into four sub-research questions:

\begin{enumerate}
    \item \textbf{Correlation}: Is there a statistically significant correlation between a model's generalisation gap and the rule explosion (complexity versus fidelity) observed in extracted symbolic rules?
    \item \textbf{Regularisation Impact}: How does the transition from standard Adam to SAM-based optimisation affect the stability and sharpness of the model's error landscape, and the stability of its explanations?
    \item \textbf{Multi-Criteria Validation}: To what degree do SAM-optimised rules outrank baseline rules when evaluated through the PROMETHEE II multi-criteria framework?
    \item \textbf{Biological Alignment}: Do the rules extracted from a regularised model show higher enrichment in KEGG/Enrichr pathways compared to rules from an overfitted baseline?
\end{enumerate}

The hypothesis is that applying SAM regularisation to a neural network trained on GDSC genomic
data will reduce the generalisation gap by converging on flatter minima, resulting in symbolic rules
that are statistically simpler, more accurately aligned to the black-box logic of the model (higher
fidelity), and more consistently aligned with known biological pathways (higher biological validity)
than those produced by an unregularised baseline \cite{Foret2021, Li2018, Springer2024, Yedjour2024}.

% -------------------------------------------------------
% NEW CONTENT REQUIRED (COS700 Slide 21)
% -------------------------------------------------------

\subsection*{Scope and Limitations}

\textcolor{red}{
% ADD: Scope and limitations paragraph as required by proposal guidelines.
% Suggested text (edit as needed):
This study is scoped to the following boundaries. The neural network architecture is fixed to
DeepCCDS \cite{Wu2025} and is not subject to architectural optimisation. The primary dataset is
GDSC with LN IC\textsubscript{50} as the drug response metric; the study does not extend to
clinical patient data or real-time deployment contexts. The rule extraction pipeline operates on
discretised tabular genomic features and does not address image-based or multi-modal data
modalities. The regularisation intervention is limited to SAM as a replacement for the Adam
optimiser; other regularisation strategies such as dropout or L2 weight decay are not evaluated as
primary interventions. Statistical comparisons are conducted between the baseline and
SAM-regularised conditions only; the study does not benchmark against all possible optimisers.
Biological verification is limited to KEGG and Enrichr pathway databases as accessed through
GSEApy; wet-lab experimental validation is outside the scope of this project.
}